\section{Thick petal sheets}\label{sec:sheets}
\stagebox{\statusbuilt}{\src{core/src/sheet.rs}}

Each petal patch becomes a \emph{thick sheet}: a smooth mid-surface described by
an $\NU \times \NV = 11 \times 19$ grid of control points (11 across, 19 from
base to tip), plus a thickness at every control point. The same object will carry
the petal from the closed bud through the whole opening, which is what makes
the animation smooth.

\subsection{From patch to control points}
Let the petal occupy rows $0,\dots,h-1$ of its shell (petals reaching fewer than
3 rows are dropped). Rows the patch skips get a centre and span by linear
interpolation between their neighbours; centres are unwrapped around the base
angle so they never jump across $\pm\pi$. Both sequences are then smoothed twice
with the kernel $(\tfrac14, \tfrac12, \tfrac14)$.

For control row $v = 0,\dots,\NV-1$ let $f_v = v/(\NV-1)$ and let $\rho = f_v(h-1)$
be the matching (fractional) automaton row, interpolated linearly to give a
centre $\bar\theta(\rho)$ and span $S(\rho)$. The petal's outline is shaped by a
taper that narrows the base into a claw and rounds the tip,
\[
  T(f) = \Bigl(\tfrac12 + \tfrac12\min(f/0.2,\,1)\Bigr)\Bigl(0.3 + 0.7\sqrt{\max(0,\,1 - f^4)}\Bigr),
\]
and the control points of row $v$ lie on the layer's shell:
\[
  \x_{u,v} = P\!\left(a_k,\; \bar\theta(\rho) + \bigl(f_u - \tfrac12\bigr)\,0.92\,S(\rho)\,T(f_v),\; \phi\!\left(\frac{\rho + \frac12}{\NP}\right)\right),
  \qquad f_u = \frac{u}{\NU - 1}.
\]
The factor $0.92$ leaves a small gap between neighbouring petals at the start.

\subsection{Thickness, capped}
Thickness is greatest at the base and along the midline and falls toward the
edges and tip:
\[
  b(f_u, f_v) = (1 - f_v)^{1.2}\bigl(1 - 0.6\,(2f_u - 1)^2\bigr) \in [0,1],
  \qquad
  t_{u,v} = \tmax\bigl(e + (1-e)\,b\bigr).
\]
Since $b \le 1$, every thickness is at most $\tmax$, so petals always stay
thin, and since $b\ge0$ the edge never drops below $e\,\tmax$. A test checks
the cap on several seeds.

\subsection{Evaluating the surface}
The control grid is turned into a finer surface by interpolating Catmull--Rom
splines, first across each control row and then along each resulting column.
Ends are handled by reflecting the second point through the first (linear
extrapolation). Thickness is interpolated linearly on the same grid. With
$s_a$ samples per span along and $s_c$ across, the evaluated grid has
$(\NV - 1)s_a + 1$ rows and $(\NU - 1)s_c + 1$ columns. Because Catmull--Rom
splines pass through their control points, the evaluated surface passes
through every control point.

